Potency labels drawn from independent ChEMBL documents disagree by 0.31 to 0.92 pChEMBL units depending on the target, and that disagreement sets a floor under the test error of any model. Estimating the floor per task, rather than borrowing one number from the literature, is cheap and changes what a leaderboard means: on 16 of our 41 tasks a plain ECFP model scores below its own floor under a random split, which is impossible if the floor is right and the test set is independent.

It is the test set that is not independent. Document support --- the number of training molecules drawn from a test molecule's own source document --- moves test RMSE from 1.07 to 0.51 pChEMBL units, and with support held at zero the random, scaffold and temporal splits are equally hard. A measured per-document offset accounts for 57\% of the noise variance but at most half of that advantage, so what a model recovers from a familiar document is mostly the series of analogues reported alongside the test molecule, not the laboratory's calibration.

Two practical recommendations follow. First, report the distribution of document support over the test set alongside any split-based score on ChEMBL data; the document identifier needed for it is already in the database. Second, where the question is generalisation rather than interpolation, group the split by source document: it removes most of the support, puts every task back below the ceiling, and leaves the ordering of models unchanged. The noise ceiling is then best read not as a constant to divide by, but as an instrument that detects when a benchmark has stopped measuring what it claims to measure.
